\section{Limiti dei modelli linguistici in un dominio specifico}
\label{sec:problema}

Un modello linguistico di grandi dimensioni, per quanto addestrato su un corpus sterminato di testo, soffre di due limiti strutturali quando interrogato direttamente su un dominio specifico e circoscritto come quello universitario.

Il primo è la natura statica della sua conoscenza: i parametri del modello fissano ciò che esso ``sa'' al momento in cui l'addestramento si conclude, e qualunque informazione successiva a quella data come una nuova scadenza, un regolamento modificato, un importo delle tasse aggiornato, semplicemente non vi rientra. 
In \cite{zhao2023survey}, nella loro rassegna sui large language model, descrivono questo fenomeno con il termine \emph{outdatedness hallucination}: il modello genera un contenuto che era corretto al momento dell'addestramento ma non lo è più al momento dell'uso, senza alcun segnale esplicito che lo distingua da un'affermazione ancora valida, il modello stesso non ha modo di sapere che la propria conoscenza è invecchiata.

Il secondo limite riguarda l'assenza totale di accesso ai documenti specifici dell'istituzione interrogata: un LLM generico non ha mai visto i regolamenti di un ateneo, le pagine web di un singolo corso di laurea, le slide di un insegnamento particolare. Messo di fronte a una domanda che richiede questo tipo di conoscenza, non ha altra scelta che generalizzare da pattern appresi su domini affini o, più problematicamente, produrre una risposta sintatticamente corretta e sicura di sé ma priva di fondamento fattuale, un comportamento che prende il nome di \emph{hallucination}. Huang et al. \cite{huang2023survey} ne propongono una tassonomia sistematica, riconducendone le cause a limiti dei dati di addestramento (rumorosi, incompleti o distorti), a processi di allineamento che premiano risposte sicure più che corrette, e alla natura intrinsecamente stocastica del processo di
generazione del testo.

Questi due limiti si combinano in modo particolarmente dannoso nel dominio universitario, dove il costo di un errore non è astratto ma concreto e misurabile: una scadenza riportata in modo errato comporta una tassa pagata in ritardo, un requisito di ammissione frainteso comporta una domanda respinta. Kasneci et al. \cite{kasneci2023chatgpt}, discutendo opportunità e criticità dell'uso dei large language model in ambito educativo, osservano come l'assenza di garanzie sull'accuratezza fattuale sia uno degli ostacoli principali alla loro adozione diretta in contesti dove la correttezza dell'informazione fornita allo studente non è negoziabile.

La soluzione che la letteratura ha individuato per entrambi questi limiti consiste nel separare la conoscenza fattuale dal componente che genera il linguaggio: non chiedere al modello di ``ricordare'' un fatto a partire dai propri parametri, ma fornirglielo esplicitamente come contesto al momento della domanda, recuperandolo da una base documentale aggiornabile in modo indipendente dal modello stesso. È il principio alla base del Retrieval-Augmented Generation, discusso nella Sezione~\ref{sec:rag}, ed è il principio su cui si fonda l'intero sistema sviluppato in questo lavoro di tesi.

Questa non è l'unica strada percorribile: la conoscenza mancante di un LLM potrebbe, in linea di principio, essere aggiunta direttamente ai suoi parametri con un ulteriore addestramento mirato (\emph{fine-tuning}) sui documenti specifici del dominio, piuttosto che recuperata da una base esterna a ogni domanda. Ovadia et al. \cite{ovadia2024finetuning} confrontano empiricamente le due strade su compiti che richiedono conoscenza intensiva, sia per informazioni già presenti nei dati di addestramento sia per informazioni del tutto nuove, e osservano che il fine-tuning non supervisionato fatica sistematicamente ad apprendere fatti nuovi in modo affidabile --- a meno di esporre il modello a numerose riformulazioni dello stesso fatto durante l'addestramento, un costo che cresce con
ogni aggiornamento del documento sorgente --- mentre il RAG si dimostra costantemente più efficace in entrambi i casi. Per un dominio come quello universitario, dove i documenti cambiano ogni anno accademico, questo risultato ha un peso pratico diretto: un sistema a fine-tuning richiederebbe un nuovo ciclo di addestramento a ogni aggiornamento normativo, mentre un sistema RAG richiede solo di re-indicizzare il documento modificato.

Questo principio ha già trovato applicazione in sistemi conversazionali rivolti specificamente al contesto
universitario e industriale. Una presentazione più ampia di sistemi analoghi è presentata nella Sezione~\ref{sec:sistemi-simili}.